\label{tab:bank_shocks_desc}
\begin{threeparttable}
\begin{tabular}{lccccccccc}
\toprule
% \multicolumn{7}{c}{\textbf{Panel A. Overall}} \\
% \midrule
%  & & & & & \textbf{Raw} & \textbf{Residualized} \\
% \midrule
% \multicolumn{5}{l}{Mean ($\mu$)} & -0.18 & 0.00 \\
% \multicolumn{5}{l}{Standard deviation ($\sigma$)} & 0.22 & 0.21 \\
% \multicolumn{5}{l}{Interquartile range (IQR)} & 0.00 & 0.00 \\
% \midrule
% \multicolumn{5}{l}{Effective sample size (ESS)} & \multicolumn{2}{c}{1.20} \\
% \multicolumn{5}{l}{Largest $s_{nt}$ weight} & \multicolumn{2}{c}{0.91} \\
% \multicolumn{5}{l}{No. of bank-year shocks} & \multicolumn{2}{c}{137} \\
% \multicolumn{5}{l}{No. of banks} & \multicolumn{2}{c}{80} \\
% % \multicolumn{4}{l}{ESS / No. of banks} & 0.042 \\
% \midrule
% \multicolumn{7}{c}{\textbf{Panel B. Yearly}} \\
% \midrule
\textbf{Year} & Mean & Std. deviation & IQR &  \textbf{No. banks} & \textbf{ESS} & \textbf{Largest w} \\
\midrule
% 2006 & -1.29 & 1.40 & 2.60 & 20 & 2.95 & 0.554 \\
% 2007 &  0.03 & 0.89 & 0.44 & 23 & 8.05 & 0.173 \\
% 2008 & -0.10 & 0.41 & 0.32 & 28 & 8.26 & 0.188 \\
2009 & 0.15 & 1.13 & 0.56 & 27 & 8.21 & 0.18 \\
% 2010 & -0.50 & 0.90 & 0.14 & 58 & 8.49 & 0.179 \\
% 2011 & -0.18 & 0.97 & 2.04 & 61 & 7.11 & 0.249 \\
% 2012 & -0.21 & 0.70 & 0.37 & 64 & 7.33 & 0.211 \\
% 2013 & -0.15 & 0.43 & 0.40 & 58 & 7.07 & 0.265 \\
2014 & 0.08 & 1.37 & 0.61 & 50 & 7.59 & 0.28 \\
% 2015 & -0.05 & 0.45 & 0.29 & 52 & 3.48 & 0.504 \\
% 2016 &  0.08 & 0.23 & 0.16 & 60 & 3.85 & 0.444 \\
% 2017 & -2.09 & 1.45 & 2.19 & 59 & 3.49 & 0.460 \\
% 2018 & -0.36 & 0.76 & 0.60 & 74 & 8.82 & 0.237 \\
2019 &  -0.05 & 0.36 & 0.52 & 72 & 9.74 & 0.20 \\
% 2020 & -0.10 & 0.25 & 0.33 & 77 & 8.24 & 0.287 \\
% 2021 &  0.16 & 0.62 & 0.27 & 77 & 9.77 & 0.225 \\
% 2022 & -0.61 & 1.54 & 0.52 & 76 & 10.19 & 0.190 \\
% 2023 &  0.02 & 0.31 & 0.00 & 72 & 1.17 & 0.924 \\
2024 & -0.07 & 1.00 & 0.71 & 67 & 1.89 & 0.72 \\
% 2025 & -0.25 & 0.13 & 0.00 & 70 & 1.11 & 0.948 \\
\bottomrule
\end{tabular}
\begin{tablenotes}[flushleft]
\footnote
\item This table summarizes the distribution of bank-level credit shocks $\hat{\gamma}_{bt}^{std}$ across banks $b$ and years $t$, when excluding markets with high concentration of bank loans (Mexico city, Estado de Mexico, Leon, Monterrey, and Queretaro). The shock corresponds to the standardized (overall the years) bank fixed effects. 

The effective sample size (ESS) is computed as $1/\sum w_{bt}^2$, and the largest weight is the maximum exposure share.

The Inter-Quartile Range (IQR) is defined as the difference between the 25 and the 75 percentiles.
%\tnote{a} 
% \item[a] Measured as $1/\text{HHI}$ of the relevant shock weights.
\end{tablenotes}
\end{threeparttable}
